\subsection{LiDAR-Messsystem RIEGL LMS VQ780II}

Die vom GeoSN bereitgestellten LiDAR-Daten wurden mit dem \textit{ALS RIEGL LMS VQ780II }erfasst. Dieses System ist speziell für die Vermessung von Gelände- und Vegetationsstrukturen konzipiert und weist sowohl eine Präzision (Wiederholungsgenauigkeit) als auch eine absolute Genauigkeit von jeweils \textit{20 mm} auf. Im Gegensatz zu \textit{Solid-State-LiDAR-Systemen}, die rein auf Halbleiterelektronik basieren, nutzt der RIEGL-Scanner ein mechanisches Scanverfahren mittels eines kontinuierlichen Polygonspiegelrads. Diese lenkt den ausgesendeten Laserstrahl zur zeilenweisen Erfassung ab, wodurch eine zeilenweise Erfassung der Oberfläche ermöglicht wird. Die Rotation erzeugt exakte, parallele Scanlinien und führt zu einem homogenen Punktmuster über das gesamte \textit{FoV} von \textit{60°}. Technisch erzielt das System eine Laserpuls-Wiederholungsrate von bis zu \textit{2,0 MHz}, woraus bis zu \textit{1,33 Millionen} Messungen pro Sekunde auf dem Boden resultieren. Da der Scanner bis zu \textit{35} Laserpulse gleichzeitig in der Luft verarbeiten kann (Multiple-Time-Around), ist eine effiziente Datenerfassung aus großen Flughöhen mit einer maximalen Messreichweite von bis zu \textit{6.800 m} möglich. Eine integrierte, effektive Unterdrückung atmosphärischer Störungen minimiert dabei den anschließendenFilteraufwand bei der Generierung bereinigter Punktwolken. Die Leistungsfähigkeit des Scanners wird jedoch durch starkes Umgebungslicht gemindert \cite{zotero-item-48}.